\chapter{Sixth-Order Resolution: What the Richer Algebra Can See}

\section{New Invariant Information}

The sixth-order invariant space for $V_5$ has dimension six:

\[
\dim\operatorname{Inv}_6(V_5)=6.
\]

The unique quadratic invariant can multiply each of the two quartic invariants, producing a two-dimensional inherited subspace $I_2\operatorname{Inv}_4$. Thus

\[
\dim(I_2\operatorname{Inv}_4)=2,
\]

leaving

\[
6-2=4
\]

genuinely new degree-six invariant directions. The dimension count already proves that sixth order contains relational information not recoverable merely by multiplying lower-order invariants. The remaining question is whether any of that new information is sensitive to the two physical shape directions left flat by quartic order.

An exact basis for the nonrotational quartic-flat tangent space at $A_\star$ can be chosen as

\[
U=-\frac{\sqrt{14}}{14}Y^c_{52}+Y^c_{53},
\]

\[
V=\frac{\sqrt{14}}{14}Y^s_{52}+Y^s_{53}.
\]

These satisfy

\[
\|U\|^2=\|V\|^2=\frac{15}{14},
\qquad
\langle U,V\rangle=0.
\]

They span the orthogonal complement of the rotational orbit within the five-dimensional quartic tangent kernel. In this sense the shape-flat plane is not chosen arbitrarily after the fact; it is the residual physical tangent subspace once normalization, the quartic constraint, and infinitesimal rotations have been accounted for.

\section{An Explicit Sextic Invariant}

Define the degree-four covariant

\[
B_4(A)=P_4(A\otimes A),
\]

then couple it once more with $A$ and project to degree eight:

\[
C_8(A)=P_8(B_4(A)\otimes A).
\]

The squared norm

\[
J_{4,8}(A)=\|C_8(A)\|^2
\]

is a scalar invariant of total degree six. Equivalently,

\[
J_{4,8}(A)=\left\|P_8\left(P_4(A\otimes A)\otimes A\right)\right\|^2.
\]

Because the quartic minimum set is constrained by both normalization and $B_2(A)=0$, second derivatives along $U$ and $V$ must be computed intrinsically rather than by simply taking straight lines in the ambient representation space. The corrected paths include second-order components chosen so that the unit-sphere condition and the quartic constraint continue to hold through second order. With that correction included, the first derivatives vanish:

\[
dJ_{4,8}(U)=dJ_{4,8}(V)=0.
\]

The intrinsic second derivatives are

\[
d^2J_{4,8}(U,U)=\frac{72}{143},
\]

\[
d^2J_{4,8}(V,V)=\frac{72}{143},
\]

and

\[
d^2J_{4,8}(U,V)=0.
\]

Normalize the shape directions by

\[
e_1=\sqrt{\frac{14}{15}}U,
\qquad
e_2=\sqrt{\frac{14}{15}}V.
\]

The intrinsic shape Hessian becomes

\[
H_{\mathrm{shape}}(J_{4,8})=\frac{336}{715}I_2\succ0.
\]

The result is exact at $A_\star$: this sextic invariant supplies positive curvature in both shape directions that quartic order leaves flat.

\begin{proposition}[Sixth-order resolving capability]
The $SO(3)$-invariant degree-six algebra on $V_5$ contains invariant structure whose intrinsic Hessian is positive definite on the two-dimensional nonrotational quartic-flat tangent space at $A_\star$. Hence sixth-order invariant theory is algebraically capable of resolving both quartic shape directions.
\end{proposition}

\section{Capability Is Not Yet PDE Realization}

The result does not yet establish the complete sixth-order reduced energy of the original nonlinear PDE. That calculation requires the center-manifold or slaving expansion through the required orders, including contributions from higher harmonics, the terms conventionally denoted $w^{(3)}$ and $w^{(4)}$, and the relevant Clebsch--Gordan recouplings. The PDE-derived coefficient multiplying $J_{4,8}$ or any equivalent primitive sextic invariant must be computed rather than assumed. A positive-definite invariant exists; the specific dynamics must still be shown to use it with the sign and magnitude required for stabilization.

This boundary is important because it illustrates the book's central distinction in miniature. The invariant algebra establishes admissibility and resolving capacity. The PDE decides occupancy of that algebraic possibility within the actual reduced energy. One should not infer the latter from the former.

\section{A General Resolution Principle}

The $\ell=5$ calculation suggests a general theorem about descriptive algebras. Let $\mathcal I_p$ be a family of invariants available at descriptive order $p$, and define

\[
x\sim_p y
\quad\Longleftrightarrow\quad
f(x)=f(y)\ \text{for all}\ f\in\mathcal I_p.
\]

If

\[
\mathcal I_p\subseteq\mathcal I_q,
\]

then

\[
[x]_q\subseteq[x]_p.
\]

A richer invariant language can leave an equivalence class unchanged or refine it. If some $g\in\mathcal I_q$ distinguishes two states that every element of $\mathcal I_p$ identifies, then the inclusion is strict. The conclusion is therefore

\begin{principle}[Relational resolution]
Indistinguishability under one restricted descriptive algebra does not imply identity under every richer descriptive algebra.
\end{principle}

The principle does not imply an infinite regress in which every apparent identity hides an unknowable difference. For compact group actions, sufficiently rich invariant algebras can separate distinct orbits under standard hypotheses. The point is instead that claims of completeness must be earned relative to a demonstrated separating structure. A finite truncation should not be treated as complete merely because no distinction is visible inside that truncation.

This becomes important when physical sufficiency is discussed later. The $\ell=5$ result does not provide an analogue of consciousness, but it does provide a warning against confusing the resolving power of the current physical description with the total structure of the physical state. A physicalist test must therefore enrich the physical description systematically before interpreting a residual as evidence of nonphysical structure.
